\begin{table}[htbp]
\centering
\caption{Marijuana arrests in non-legal border counties around a neighboring state's legalization}
\label{tab:arrestdid}
\footnotesize
\begin{tabular}{lccccc}
\toprule
Legalizing neighbor & Treated counties & Control counties & $N$ & DID coef. (\%) & $p$ \\
\midrule
Colorado (2014-01) & 10 & 175 & 12,395 & $+0.688$ ($+98.9\%$) & $0.003$ \\
Arizona (2021-01)  & 4  & 20  & 1,608  & $+0.632$ ($+88.1\%$) & $0.148$ \\
New Mexico (2022-04) & 33 & 80 & 7,571  & $-0.075$ ($-7.2\%$)  & $0.594$ \\
\midrule
Pooled (stacked) & 47 & 175\textsuperscript{a} & 21,574 & $+0.186$ ($+20.4\%$) & $0.148$ \\
\bottomrule
\end{tabular}
\begin{minipage}{0.98\textwidth}
\vspace{0.3em}
\footnotesize
\textit{Notes:} Dependent variable is $\log(1+\text{marijuana arrests})$ per county-quarter (FBI offense
codes 158 + 153). Each row is a separate stacked-DID estimate of $\pi$ from equation~\ref{eq:arrestdid},
restricted to that cohort's own treated and control counties; the pooled row stacks all three cohorts
together. County and cohort-by-quarter fixed effects; county-clustered standard errors.
\textsuperscript{a}Unique control counties across all three stacks -- a Texas or Utah interior county can
serve as a control for more than one cohort, so this is less than the sum of the per-cohort control columns.
\end{minipage}
\end{table}
